\documentclass[../../main.tex]{subfiles}
\graphicspath{{\subfix{./image/}}{\subfix{../../image/}}}
\begin{document}
% Pagination repair: keep the new section mark off the preceding section's last page.
\clearpage
% SOURCE p=158; printed=143; object=ch04-cube-embeddings-title; region=middle
\section{嵌入立方体}\label{section:ch04-cube-embeddings}
% SOURCE p=158; printed=143; object=ch04-cube-embeddings-intro; region=middle
现在介绍把拓扑空间嵌入区间乘积的一种方法，并讨论若干应用。（本书其他部分不会使用这些结果。）本节始终用 $I$ 表示单位区间 $[0,1]$。
% SOURCE p=158; printed=143; object=ch04-cube-definition; region=middle
\begin{definition}[立方体]\label{definition:ch04-cube}
若 $A$ 是任意非空集合，则称乘积空间 $I^A$ 为一个\textbf{立方体}。
\end{definition}
% SOURCE p=158; printed=143; object=ch04-points-closed-separation; region=middle
\begin{definition}[分离点与闭集的函数族]\label{definition:ch04-points-closed-separation}
设 $X$ 为拓扑空间，$\mathcal F\subseteq C(X,I)$。若对每个闭集 $E\subseteq X$ 及每个 $x\in E^c$，存在 $f\in\mathcal F$ 使 $f(x)\notin\overline{f(E)}$，则称 $\mathcal F$ \textbf{分离点与闭集}。
\end{definition}
% SOURCE p=158; printed=143; object=ch04-strong-points-closed-separation; region=middle
\begin{proposition}\label{proposition:ch04-strong-points-closed-separation}
若 $\mathcal F$ 分离点与闭集，则存在另一个函数族 $\mathcal G\subseteq C(X,I)$，满足稍强的性质：对每个闭集 $E\subseteq X$ 及每个 $x\in E^c$，存在 $g\in\mathcal G$ 使 $g(x)=1$ 且 $g=0$ 于 $E$。因此，一个 $T_1$ 空间具有分离点与闭集的函数族，当且仅当它完全正则。
\end{proposition}
% SOURCE p=158; printed=143; object=ch04-proof-strong-points-closed-separation; region=middle
\begin{proof}
若 $f\in\mathcal F$ 满足 $f(x)\notin\overline{f(E)}$，取 $g=\phi\circ f$，其中 $\phi\in C(I,I)$、$\phi(f(x))=1$，且 $\phi=0$ 于 $\overline{f(E)}$。
\end{proof}
% SUPPLEMENT: id=v14-ch04-cube-separation-details; parent=ch04-strong-points-closed-separation; kind=proof-detail+motivation+hypotheses; origin=independent-derivation
\begin{proof}
\textbf{补写细节（重构新增）。}闭区间 $I$ 是正规空间，$\{f(x)\}$ 与 $\overline{f(E)}$ 是不交闭集，故 Urysohn 引理给出所需 $\phi$；若 $E=\varnothing$，也可直接取 $\phi=1$。取
\begin{align*}
\mathcal G=\{\phi\circ f:f\in\mathcal F,\ \phi\in C(I,I)\},
\end{align*}
便得到一个同时适用于全部点与闭集的族。上述强分离性质与 $T_1$ 合在一起，正是完全正则性的定义；反向若完全正则，取 $\mathcal F=C(X,I)$，由值 $1$ 与 $0$ 的分离便满足较弱的定义。
\end{proof}
\begin{remark}
\textbf{定义的动机与条件说明（重构新增）。}立方体使用乘积拓扑，每个基本邻域只约束有限多个坐标；任意多个紧区间的积仍紧 Hausdorff，因而适合作为统一的容纳空间。这里的“立方体”不局限于有限维。原书要求指标非空，以免讨论空积；若允许空积，它是单点空间。

分离点与闭集要求 $f(x)$ 连 $f(E)$ 的极限值也避开，这样才能在值域中找到一个开邻域，并把它拉回为避开 $E$ 的邻域。只要求 $f(x)\notin f(E)$ 不能保证这个步骤。函数族中不一定已有取值恰为 $0,1$ 的分离函数，所以原命题允许改用复合得到的另一族。$E$ 闭、$x\notin E$ 是完全正则性检验的对象范围；$T_1$ 是本书完全正则定义中的独立分离条件。
\end{remark}
% END SUPPLEMENT: id=v14-ch04-cube-separation-details
% SOURCE p=158; printed=143; object=ch04-associated-cube-map; region=middle-lower
\begin{definition}[函数族的关联映射]\label{definition:ch04-associated-cube-map}
每个非空函数族 $\mathcal F\subseteq C(X,I)$ 都典范地诱导映射 $e:X\to I^{\mathcal F}$，由
\begin{align*}
\pi_f(e(x))=f(x)\qquad(f\in\mathcal F)
\end{align*}
定义，其中 $\pi_f:I^{\mathcal F}\to I$ 为坐标映射。称 $e$ 为与 $\mathcal F$ \textbf{关联的映射}。
\end{definition}
% SOURCE p=158; printed=143; object=ch04-associated-map-generalization; region=lower
显然，这个构造可推广到 $I$ 以外的目标空间，见\cref{exercise:ch04-product-universal}。
% SOURCE p=158; printed=143; object=ch04-associated-embedding; number=4.53; region=lower
\begin{proposition}\label{proposition:ch04-associated-embedding}
设 $X$ 为拓扑空间，$\mathcal F\subseteq C(X,I)$，$e:X\to I^{\mathcal F}$ 是与 $\mathcal F$ 关联的映射。则：
\begin{enumerate}
\item\label{item:ch04-associated-continuous} $e$ 连续。
\item\label{item:ch04-associated-injective} 若 $\mathcal F$ 分离点，则 $e$ 单射。
\item\label{item:ch04-associated-embedding} 若 $X$ 为 $T_1$ 空间且 $\mathcal F$ 分离点与闭集，则 $e$ 是嵌入。
\end{enumerate}
\end{proposition}
% SOURCE p=158; printed=143; object=ch04-proof-associated-embedding; region=lower
\begin{proof}
\subcref{proposition:ch04-associated-embedding}{item:ch04-associated-continuous}由\cref{proposition:ch04-product-map-continuity}得出，\subcref{proposition:ch04-associated-embedding}{item:ch04-associated-injective}显然。接着注意：若 $\mathcal F$ 分离点与闭集，且 $X$ 为 $T_1$，则由\subcref{proposition:ch04-associated-embedding}{item:ch04-associated-injective}和\cref{proposition:ch04-t1-singletons}知 $e$ 单射。为证明逆映射连续，设 $U$ 在 $X$ 中开。若 $x\in U$，选择 $f\in\mathcal F$ 使 $f(x)\notin\overline{f(U^c)}$，并令
\begin{align*}
V=\pi_f^{-1}\bigl(I\setminus\overline{f(U^c)}\bigr)
=\{p\in I^{\mathcal F}:\pi_f(p)\notin\overline{f(U^c)}\}.
\end{align*}
则 $V$ 在 $I^{\mathcal F}$ 中开，且 $e(x)\in V\cap e(X)\subseteq e(U)$。所以对每个 $x\in U$，$e(U)$ 都是 $e(x)$ 在 $e(X)$ 中的邻域，因而 $e(U)$ 在 $e(X)$ 中开。这说明 $e^{-1}$ 连续。
\end{proof}
% SUPPLEMENT: id=v14-ch04-associated-embedding-details; parent=ch04-associated-embedding; kind=proof-detail+motivation+hypotheses; origin=independent-derivation
\begin{proof}
\textbf{单射与逆连续性的补写（重构新增）。}每个坐标复合 $\pi_f\circ e=f$ 连续，故 $e$ 连续。若 $x\ne y$，分点函数使至少一个坐标值不同，故 $e(x)\ne e(y)$。在第三项中，$T_1$ 使 $\{y\}$ 闭，对点 $x$ 与该闭集应用分离性质便得到这样的函数。

原证中的包含关系也可逐点验证：若 $e(z)\in V$，则 $f(z)\notin\overline{f(U^c)}$，所以 $z\notin U^c$，即 $z\in U$。把随 $x\in U$ 选择的 $V$ 与 $e(X)$ 相交后取并，恰得 $e(U)$；因此所有开集的像在像空间中开，逆映射连续。空 $U$ 的像为空，自动开。
\end{proof}
\begin{remark}
\textbf{动机与逐条件说明（重构新增）。}关联映射把一个点变成其全部“函数读数”。函数连续使读数随点连续变化，分点性质使不同点的读数不同，而分离点与闭集使读数还足以恢复原有邻域，故从连续单射升级为嵌入。$T_1$ 在此用于把另一个点视作闭集；若完全删去分离假设，两个点的不可分拓扑空间只有常值连续函数，却仍可用常函数族分离点与其唯一可能的不含该点的闭集 $\varnothing$，关联映射并不单射。非空函数族沿用前一定义；需要时加入常函数，不会损害任何分离性质。
\end{remark}
% END SUPPLEMENT: id=v14-ch04-associated-embedding-details
% SOURCE p=159; printed=144; object=ch04-compact-closed-cube; number=4.54; region=top
\begin{corollary}\label{corollary:ch04-compact-closed-cube}
每个紧 Hausdorff 空间都同胚于某个立方体的闭子集。
\end{corollary}
% SOURCE p=159; printed=144; object=ch04-proof-compact-closed-cube; region=top
\begin{proof}
由\cref{proposition:ch04-compact-hausdorff-normal}及 Urysohn 引理，可取 $\mathcal F=C(X,I)$。
\end{proof}
% SUPPLEMENT: id=v14-ch04-compact-closed-cube-proof; parent=ch04-compact-closed-cube; kind=proof-detail+motivation+hypotheses; origin=independent-derivation
\begin{proof}
\textbf{补写完整链条（重构新增）。}紧 Hausdorff 空间正规并为 $T_1$。任意闭集 $E$ 与外部点 $x$ 可由 Urysohn 函数取值 $0,1$ 分离，故 $C(X,I)$ 分离点与闭集。由\cref{proposition:ch04-associated-embedding}，关联映射 $e$ 为嵌入。连续像 $e(X)$ 紧，而立方体是 Hausdorff 空间，所以 $e(X)$ 闭。若 $X$ 为空，则 $C(X,I)$ 仍含唯一空函数，空像为立方体的闭子集，论证兼容。
\end{proof}
\begin{remark}
\textbf{条件与动机（重构新增）。}这里不仅把紧空间放进统一模型，还保留“闭子集”这个紧性标志。Hausdorff 与紧性先保证足够多分离函数；紧性随后使嵌入像紧，目标的 Hausdorff 性再使该像闭。一般完全正则空间只保证嵌入像是子集，不保证闭。
\end{remark}
% END SUPPLEMENT: id=v14-ch04-compact-closed-cube-proof
% SOURCE p=159; printed=144; object=ch04-completely-regular-compact-subspace; number=4.55; region=upper
\begin{corollary}\label{corollary:ch04-completely-regular-compact-subspace}
拓扑空间完全正则，当且仅当它同胚于某个紧 Hausdorff 空间的子集。
\end{corollary}
% SOURCE p=159; printed=144; object=ch04-proof-completely-regular-compact-subspace; region=upper
\begin{proof}
在\cref{proposition:ch04-associated-embedding}中取 $\mathcal F=C(X,I)$，得到“仅当”方向；反向留给读者，见\cref{exercise:ch04-completely-regular-subspace}。
\end{proof}
% SUPPLEMENT: id=v14-ch04-completely-regular-subspace-proof; parent=ch04-completely-regular-compact-subspace; kind=omitted-proof+motivation+hypotheses; origin=independent-derivation
\begin{proof}
\textbf{反向与继承性的补写（重构新增）。}设 $Z$ 完全正则，$S\subseteq Z$ 取相对拓扑。单点在 $Z$ 中闭，所以也在 $S$ 中闭，$S$ 为 $T_1$。若 $E\subseteq S$ 相对闭且 $x\in S\setminus E$，可写 $E=S\cap F$，其中 $F$ 在 $Z$ 中闭。由 $x\in S$ 且 $x\notin E$ 知 $x\notin F$。取 $g\in C(Z,I)$ 使 $g(x)=1$、$g|_F=0$，其限制 $g|_S$ 就分离 $x$ 与 $E$，故 $S$ 完全正则。

紧 Hausdorff 空间由正规性与 Urysohn 引理完全正则，于是其每个子空间完全正则。完全正则性在同胚下保持：单点闭性由同胚传递，分离函数与同胚复合即可。因此得到原推论的反向。正向的立方体本身由 Tychonoff 定理紧 Hausdorff，补齐了原证的目标空间条件。
\end{proof}
\begin{remark}
\textbf{动机与条件说明（重构新增）。}完全正则性恰好描述哪些空间可以成为紧 Hausdorff 空间的一部分。这里的子集使用相对拓扑，且不要求闭；若要求闭，紧性会被继承，从而排除非紧的完全正则空间。继承性证明中的环境闭集 $F$ 不必等于 $E$，这是处理非闭子空间时的关键。
\end{remark}
% END SUPPLEMENT: id=v14-ch04-completely-regular-subspace-proof
% SOURCE p=159; printed=144; object=ch04-compactification-definition; region=upper-middle
\begin{definition}[紧化]\label{definition:ch04-compactification}
拓扑空间 $X$ 的一个\textbf{紧化}是二元组 $(Y,\phi)$，其中 $Y$ 为紧 Hausdorff 空间，而 $\phi$ 是从 $X$ 到 $Y$ 的某个稠密子集的同胚。常把 $X$ 与其像 $\phi(X)\subseteq Y$ 等同，此时简称“$X$ 的紧化 $Y$”。
\end{definition}
% SOURCE p=159; printed=144; object=ch04-compactification-examples; region=middle
\begin{example}\label{example:ch04-compactification-examples}
$([-1,1],\tanh)$ 是 $\mathbb R$ 的一个紧化；LCH 空间 $X$ 的单点紧化 $(X^*,i)$ 也是这里意义下的紧化，其中 $i:X\to X^*$ 是包含映射。
\end{example}
% SOURCE p=159; printed=144; object=ch04-associated-compactification-construction; region=middle
设 $X$ 完全正则。根据\cref{proposition:ch04-associated-embedding}，若 $\mathcal F\subseteq C(X,I)$ 分离点与闭集，$e:X\to I^{\mathcal F}$ 为关联嵌入，而 $Y$ 为 $e(X)$ 在 $I^{\mathcal F}$ 中的闭包，则 $(Y,e)$ 是 $X$ 的一个紧化。把 $X$ 与 $e(X)$ 等同后，每个 $f\in\mathcal F$ 都有到 $Y$ 的连续延拓；由于 $X$ 在 $Y$ 中稠密，延拓唯一。事实上，这种等同把 $f$ 变成坐标映射 $\pi_f|_{e(X)}$，它延拓为 $\pi_f|_Y$。

此外，若 $f,g$ 是 $X$ 上能连续延拓到 $Y$ 的有界连续函数，则 $f+g$ 和 $fg$ 显然也能延拓。若 $\{f_n\}$ 是 $X$ 上能连续延拓到 $Y$ 的函数构成的一致收敛序列，则因 $X$ 在 $Y$ 中稠密，其延拓在 $Y$ 上也一致收敛，所以 $f=\lim f_n$ 也能连续延拓。我们已证明：
% SOURCE p=159; printed=144; object=ch04-extension-generated-algebra; number=4.56; region=middle-lower
\begin{proposition}\label{proposition:ch04-extension-generated-algebra}
设 $\mathcal F\subseteq C(X,I)$ 分离点与闭集。令 $(Y,e)$ 为与 $\mathcal F$ 关联的 $X$ 的紧化，$\mathcal A$ 为 $BC(X)$ 中包含 $\mathcal F$ 的最小闭子代数。则每个 $f\in\mathcal A$ 都有到 $Y$ 的连续延拓。
\end{proposition}
% SOURCE p=159; printed=144; object=ch04-extension-converse-reference; region=lower
这个结果有一个逆命题，见\cref{exercise:ch04-compactification-algebras}。
% SUPPLEMENT: id=v14-ch04-compactification-extension-details; parent=ch04-extension-generated-algebra; kind=proof-detail+motivation+hypotheses+source-correction; origin=independent-derivation
\begin{proof}
\textbf{闭代数与一致延拓的补写（重构新增）。}沿用前段的完全正则假设，$e$ 为嵌入，$Y$ 是紧立方体的闭子集，故紧 Hausdorff；$e(X)$ 按构造在 $Y$ 中稠密。因此确实得到紧化。

记 $R:C(Y)\to BC(X)$ 为 $Rh=h\circ e$。对任意 $h\in C(Y)$，有
\begin{align*}
\|h\|_{C(Y)}=\sup_{x\in X}|h(e(x))|=\|Rh\|_{BC(X)}.
\end{align*}
其中“$\ge$”直接；反向若 $|h(y)|$ 大于右侧，由连续性会有 $y$ 的非空开邻域，其上仍大于右侧，稠密性使该邻域遇到 $e(X)$，矛盾。空 $X$ 时稠密性迫使 $Y$ 为空，两边按零空间范数为零。

$R$ 保持复线性、乘法和单位，是等距映射。若 $Rh_n$ 在 $BC(X)$ 中一致趋于 $f$，等距性使 $h_n$ 在完备空间 $C(Y)$ 中为 Cauchy 列，故一致趋于 $h\in C(Y)$；连续限制给 $Rh=f$。因此像 $R(C(Y))$ 是闭子代数。各坐标限制 $\pi_f|_Y$ 说明它包含 $\mathcal F$，故也包含由 $\mathcal F$ 生成的最小闭子代数 $\mathcal A$，命题得证。延拓的唯一性来自两个到 Hausdorff 空间 $\mathbb C$ 的连续映射在稠密子集相等。
\end{proof}
\begin{remark}
\textbf{定义的动机与逐条件说明（重构新增）。}紧化把非紧空间嵌入紧环境以便使用紧性，但嵌入要求保留原拓扑，稠密性要求新空间完全由旧空间的极限行为决定。后者使连续延拓唯一，也使限制映射保持一致范数；紧 Hausdorff 则提供紧像闭性和连续函数分离。

$\tanh$ 严格递增、连续且把 $\mathbb R$ 映到 $(-1,1)$，逆映射 $t\mapsto\frac12\log\frac{1+t}{1-t}$ 连续；$(-1,1)$ 在 $[-1,1]$ 中稠密，所以例子成立。\textbf{单点紧化例子须沿用非紧条件：}若 $X$ 已紧，通常添加的点 $\infty$ 是孤立点，$X$ 在 $X^*$ 中不稠密，故不符合这里的紧化定义。对非紧 LCH 空间，该点的每个邻域都遇到 $X$，因此符合定义。

“最小闭子代数”指所有包含该族的闭子代数之交，存在性由整个 $BC(X)$ 保证；闭性来自一致范数。此命题只说这些函数可延拓，并没有断言它们已穷尽全部可延拓函数。原书随后习题中的逆向等号需补上常函数 $1$，详见\subcref{exercise:ch04-compactification-algebras}{item:ch04-compactification-generated-algebra}的勘误。
\end{remark}
% END SUPPLEMENT: id=v14-ch04-compactification-extension-details
% SOURCE p=159; printed=144; object=ch04-stone-cech-definition; region=lower
\begin{definition}[Stone--Čech 紧化]\label{definition:ch04-stone-cech}
若 $X$ 是完全正则空间，与 $\mathcal F=C(X,I)$ 关联的紧化称为 $X$ 的 \textbf{Stone--Čech 紧化}，记作 $(\beta X,e)$；若把 $X$ 与 $e(X)$ 等同，也简称 $\beta X$。
\end{definition}
% SOURCE p=159; printed=144; object=ch04-stone-cech-all-bounded-extend; region=lower
每个 $f\in BC(X)$ 都连续延拓到 $\beta X$；事实上，还有更一般的结果。
% SUPPLEMENT: id=v14-ch04-stone-cech-definition-details; parent=ch04-stone-cech-all-bounded-extend; kind=omitted-proof+motivation+hypotheses; origin=independent-derivation
\begin{proof}
\textbf{全部有界函数延拓的补写（重构新增）。}对有界连续实函数 $u$，取 $M>\|u\|_u$，则 $v=(u+M)/(2M)\in C(X,I)$。定义中的坐标函数使 $v$ 有连续延拓 $\widetilde v$，因而 $\widetilde u=2M\widetilde v-M$ 延拓 $u$。对复函数 $f$，分别延拓 $\operatorname{Re}f$ 与 $\operatorname{Im}f$，再相加得到延拓。唯一性仍由稠密性给出，不依赖所选 $M$。
\end{proof}
\begin{remark}
\textbf{动机与条件说明（重构新增）。}选择全部 $[0,1]$ 值连续函数作为坐标，使所有有界连续函数都能从这些坐标和常数恢复，因此新边界记录了最多的有界连续函数信息。完全正则保证原空间被嵌入；有界性不可去掉，因为紧空间上每个连续实或复函数必有界，无界函数不可能是其限制。$\beta X$ 的定义取嵌入像的闭包，并非整个庞大的立方体。
\end{remark}
% END SUPPLEMENT: id=v14-ch04-stone-cech-definition-details
% SOURCE p=159; printed=144; object=ch04-stone-cech-universal; number=4.57; region=bottom
\begin{theorem}\label{theorem:ch04-stone-cech-universal}
若 $X$ 为完全正则空间，$Y$ 为紧 Hausdorff 空间，且 $\phi\in C(X,Y)$，则 $\phi$ 有到 $\beta X$ 的唯一连续延拓 $\widetilde\phi$；即存在唯一的 $\widetilde\phi\in C(\beta X,Y)$ 使 $\widetilde\phi\circ e=\phi$。若 $(Y,\phi)$ 是 $X$ 的紧化，则 $\widetilde\phi$ 满射；若此外每个 $f\in BC(X)$ 都连续延拓到 $Y$（即对某个 $g\in C(Y)$ 有 $f=g\circ\phi$），则 $\widetilde\phi$ 为同胚。
\end{theorem}
% SOURCE p=159,160; printed=144,145; object=ch04-proof-stone-cech-universal; region=bottom-top
\begin{proof}
令 $\mathcal F=C(X,I)$、$\mathcal G=C(Y,I)$，并令 $(\beta Y,i)$ 为 $Y$ 的 Stone--Čech 紧化。也就是说，$i:Y\to I^{\mathcal G}$ 是与 $\mathcal G$ 关联的嵌入，$\beta Y=i(Y)$；因为 $Y$ 紧，$\beta Y$ 同胚于 $Y$。给定 $\phi\in C(X,Y)$，定义 $\Phi:I^{\mathcal F}\to I^{\mathcal G}$ 为
\begin{align*}
\pi_g(\Phi(p))=\pi_{g\circ\phi}(p).
\end{align*}
由\cref{proposition:ch04-product-map-continuity}知 $\Phi$ 连续，且
\begin{align*}
\pi_g(\Phi(e(x)))=\pi_{g\circ\phi}(e(x))
=g(\phi(x))=\pi_g(i(\phi(x))).
\end{align*}
即 $\Phi\circ e=i\circ\phi$。于是 $\Phi(e(X))=i(\phi(X))\subseteq\beta Y$，从而 $\Phi(\beta X)\subseteq\overline{\beta Y}=\beta Y$。这些关系可概括为下列\hyperref[figure:ch04-stone-cech-diagram]{交换图}：
% SOURCE p=160; printed=145; object=ch04-stone-cech-diagram; region=upper-middle; unnumbered=true
\begin{figure}[H]
\centering
\phantomsection\label{figure:ch04-stone-cech-diagram}
\includegraphics[width=0.38\textwidth]{ch04-stone-cech-diagram.png}
\end{figure}
令 $\widetilde\phi=i^{-1}\circ(\Phi|_{\beta X})$，则
\begin{align*}
\widetilde\phi\circ e=i^{-1}\circ\Phi\circ e=\phi.
\end{align*}
由于 $e(X)$ 在 $\beta X$ 中稠密，$\widetilde\phi$ 的唯一性显然，第一项断言得证。若 $(Y,\phi)$ 是 $X$ 的紧化，则 $\phi(X)$ 在 $Y$ 中稠密；于是 $\widetilde\phi(\beta X)$ 在 $Y$ 中既稠密又紧，所以 $\widetilde\phi(\beta X)=Y$。最后，若每个 $f\in BC(X)$ 都可写成某个 $g\in C(Y)$ 的 $g\circ\phi$，则 $\Phi$ 单射；所以 $\widetilde\phi$ 双射，从而由\cref{proposition:ch04-compact-bijection-homeomorphism}是同胚。
\end{proof}
% SUPPLEMENT: id=v14-ch04-stone-cech-universal-details; parent=ch04-stone-cech-universal; kind=proof-detail+motivation+hypotheses; origin=independent-derivation
\begin{proof}
\textbf{闭包、唯一性与最后单射步骤的补写（重构新增）。}$g\circ\phi\in C(X,I)=\mathcal F$，所以定义 $\Phi$ 使用的坐标确实存在。若不同 $g$ 给出同一复合函数，$\Phi$ 只是在这些目标坐标重复同一个源坐标，并不妨碍映射良定义。

由于 $i(Y)$ 是 Hausdorff 立方体中的紧子集，它闭，故 $\beta Y=\overline{i(Y)}=i(Y)$。连续映射把一个集合的闭包映入其像的闭包：若 $p\in\overline{e(X)}$，$\Phi(p)$ 的任意开邻域的原像都为 $p$ 的邻域，从而遇到 $e(X)$。因此
\begin{align*}
\Phi(\beta X)\subseteq\overline{\Phi(e(X))}\subseteq i(Y).
\end{align*}
$i^{-1}$ 在 $i(Y)$ 上连续，延拓构造成立。任意两个延拓的相等点集闭，并包含稠密的 $e(X)$，所以它们相等；这里使用 $Y$ Hausdorff。若 $(Y,\phi)$ 为紧化，延拓的像紧，因 $Y$ Hausdorff 而闭，又包含稠密集 $\phi(X)$，故满射。

在最后的附加假设下，对每个 $f\in\mathcal F$ 先取 $h\in C(Y)$ 使 $f=h\circ\phi$。因为 $h$ 在稠密集 $\phi(X)$ 上取值于闭集 $I$，闭集 $h^{-1}(I)$ 包含稠密集，故等于 $Y$。所以实际上 $h\in\mathcal G$。于是每个源坐标 $f$ 都等于某个 $h\circ\phi$。若 $p,q\in I^{\mathcal F}$ 不同，选 $f$ 使 $\pi_f(p)\ne\pi_f(q)$，再选相应 $h$，就有
\begin{align*}
\pi_h(\Phi(p))=\pi_f(p)\ne\pi_f(q)=\pi_h(\Phi(q)).
\end{align*}
因此 $\Phi$ 单射，限制和与 $i^{-1}$ 复合后仍单射。已知满射，紧空间到 Hausdorff 空间的连续双射遂为同胚。若 $X$ 为空，$\beta X$ 为空，唯一延拓为空映射；紧化条件还迫使 $Y$ 为空，各结论同样成立。
\end{proof}
\begin{remark}
\textbf{动机与逐条件说明（重构新增）。}这个普遍性质把标量函数的延拓提升为向任意紧 Hausdorff 空间的映射延拓：先把目标也放进立方体，再逐坐标操作。目标紧性使其嵌入像闭，从而极限不会离开目标；Hausdorff 性保证分离函数、延拓唯一及紧像闭。映射 $\phi$ 连续，才能让每个 $g\circ\phi$ 成为合法坐标。若目标不紧，恒等映射 $\mathbb R\to\mathbb R$ 不可能连续延拓到紧空间 $\beta\mathbb R$，否则其像将是包含全部 $\mathbb R$ 的紧子集。

“$\phi(X)$ 稠密”才保证延拓满射；“全部有界连续函数均可延拓”进一步保证边界点不会被合并，从而得到单射。这里只预先要求原 $X$ 完全正则，不要求它局部紧或可度量。
\end{remark}
% END SUPPLEMENT: id=v14-ch04-stone-cech-universal-details
% SOURCE p=160; printed=145; object=ch04-largest-smallest-compactifications; region=middle
这个定理说明，$\beta X$ 是完全正则空间 $X$ 的“最大”紧化，意思是其他每个紧化都是它的连续像。尺度的另一端是：若 $X$ 局部紧，则由 Urysohn 引理，$\mathcal F=C_c(X)\cap C(X,I)$ 分离点与闭集。查看与这个 $\mathcal F$ 关联的紧化 $(Y,e)$ 的构造，可以看出 $Y$ 由 $e(X)$ 加上 $I^{\mathcal F}$ 中所有坐标都为零的一个点组成。随后很容易验证，$Y$ 同胚于\cref{section:ch04-lch-spaces}中构造的 $X$ 的单点紧化。
% SUPPLEMENT: id=v14-ch04-minimal-compactification-proof; parent=ch04-largest-smallest-compactifications; kind=omitted-proof+motivation+hypotheses+source-correction; origin=independent-derivation
\begin{remark}
\textbf{原段的条件澄清（重构新增）。}这里沿用 $X$ 完全正则的总背景，因此它 Hausdorff；再假设局部紧，就得到 LCH。关于新增全零点及单点紧化的结论还必须要求 $X$ 非紧。若 $X$ 紧，则 $1\in C_c(X)\cap C(X,I)$，其坐标在 $e(X)$ 及其闭包中恒为 $1$，全零点不在闭包中。此时关联紧化只是 $X$ 自身。以下证明分别处理两种情形。
\end{remark}
\begin{proof}
\textbf{省略验证的补写（重构新增）。}对闭集 $E$ 与 $x\notin E$，局部紧 Urysohn 引理在开集 $X\setminus E$ 内给出 $f\in C_c(X,I)$，使 $f(x)=1$ 且 $\operatorname{supp}f\subseteq X\setminus E$，故 $\mathcal F$ 分离点与闭集。

先设 $X$ 非紧。令 $j:X\to X^*$ 为单点紧化包含映射，并把每个 $f\in\mathcal F$ 在 $\infty$ 处补零，得到 $\widehat f\in C(X^*,I)$：其紧支集之外函数为零，保证在无穷远连续。定义
\begin{align*}
E:X^*\longrightarrow I^{\mathcal F},\qquad
\pi_f(E(t))=\widehat f(t).
\end{align*}
它逐坐标连续，且 $E\circ j=e$、$E(\infty)=0$。函数族分离 $X$ 内各点；对每个 $x\in X$，上述截断函数可取 $f(x)=1$，于是 $E(x)\ne0=E(\infty)$。故 $E$ 单射。$X^*$ 紧，立方体 Hausdorff，所以 $E$ 为到其紧闭像的同胚。

因 $j(X)$ 在 $X^*$ 中稠密，连续性给 $E(X^*)\subseteq\overline{e(X)}$；另一方面 $E(X^*)$ 闭且包含 $e(X)$，故 $\overline{e(X)}\subseteq E(X^*)$。两边相等，且
\begin{align*}
Y=\overline{e(X)}=E(X^*)=e(X)\cup\{0\}.
\end{align*}
这既证明只有一个新点，也证明拓扑确实是单点紧化的拓扑。

若 $X$ 紧，$e(X)$ 是立方体中的紧闭集，所以 $Y=e(X)$，与 $X$ 同胚；无需也不能增加全零点。
\end{proof}
\begin{remark}
\textbf{动机与条件作用（重构新增）。}紧支集函数在远离任意固定紧区域时都变为零，因而无法区别不同的“趋向无穷远”方式，关联紧化把这些方式收成一个点。全部有界连续函数保留更多区别，得到最大紧化。局部紧性提供紧支集分离函数及 Hausdorff 单点紧化；非紧性使新增点属于旧空间的闭包，稠密性才成立。
\end{remark}
% END SUPPLEMENT: id=v14-ch04-minimal-compactification-proof
% SOURCE p=160; printed=145; object=ch04-metrization-introduction; region=lower
作为嵌入 $e:X\to I^{\mathcal F}$ 的最后一个应用，我们对下列问题给出一个部分答案：何时拓扑空间可度量化，即其拓扑何时可由某个度量定义？空间可度量化的一个必要条件是正规，见\cref{exercise:ch04-metric-normal}。另一方面：
% SOURCE p=160; printed=145; object=ch04-urysohn-metrization; number=4.58; region=lower
\begin{theorem}[Urysohn 度量化定理]\label{theorem:ch04-urysohn-metrization}
每个第二可数的正规空间都可度量化。
\end{theorem}
% SOURCE p=160; printed=145; object=ch04-proof-urysohn-metrization; region=bottom
\begin{proof}
可度量化空间的每个子空间仍可度量化，使用同一个度量的限制即可。因此，此定理直接由\cref{proposition:ch04-associated-embedding}及以下两个事实推出；它们的证明提纲见\cref{exercise:ch04-countable-separating-family,exercise:ch04-countable-product-metric}。
\begin{enumerate}
\item 若 $X$ 正规且第二可数，则存在可数函数族 $\mathcal F\subseteq C(X,I)$ 分离点与闭集。
\item 若 $\mathcal F$ 可数，则 $I^{\mathcal F}$ 可度量化。
\end{enumerate}
\end{proof}
% SUPPLEMENT: id=v14-ch04-urysohn-metrization-details; parent=ch04-urysohn-metrization; kind=proof-detail+motivation+hypotheses; origin=independent-derivation
\begin{proof}
\textbf{把两项事实接合为度量的补写（重构新增）。}下方两道习题给出这两个事实的完整证明，均不使用本度量化定理。取所构造的可数分离族；必要时加入常函数 $1$，再将其列作序列 $(f_n)_{n\ge1}$，有限族可重复列举。关联映射为嵌入，等价地由习题中的乘积度量在 $X$ 上拉回得到
\begin{align*}
d(x,y)=\sum_{n=1}^{\infty}2^{-n}|f_n(x)-f_n(y)|.
\end{align*}
分点性保证 $d(x,y)=0$ 仅当 $x=y$；其他度量公理由各项和的性质得到。由乘积度量与乘积拓扑一致、关联映射为嵌入，$d$ 诱导的拓扑正是原拓扑。空 $X$ 取唯一空度量即可。
\end{proof}
\begin{remark}
\textbf{动机与逐条件说明（重构新增）。}第二可数性使分离函数可以只选可数多个，进而用收敛的加权级数把全部函数读数压缩成一个距离。正规性通过 Urysohn 引理制造所需函数；本书“正规”包含 $T_1$，所以同时保证分点与度量的正定性。若只用不含 $T_1$ 的另一种正规约定，两点不可分拓扑虽有有限基且可真空分离不交闭集，却不可能由度量给出。

本定理是充分条件，不把第二可数性宣称为所有可度量空间的必要条件。例如不可数集合取离散度量即可度量，但每个单点必须出现在拓扑基中，故不第二可数。
\end{remark}
% END SUPPLEMENT: id=v14-ch04-urysohn-metrization-details
\Needspace{20\baselineskip}
% SOURCE p=161; printed=146; object=ch04-cube-exercises-heading; region=top
\subsection*{习题}
% SOURCE p=161; printed=146; object=ch04-ex72-completely-regular-subspace; number=72; region=top
\begin{exercise}\label{exercise:ch04-completely-regular-subspace}
完全正则空间的每个子集，在相对拓扑下都完全正则。
\end{exercise}
% SUPPLEMENT: id=v14-ch04-ex72-solution; parent=ch04-ex72-completely-regular-subspace; kind=solution+motivation+hypotheses; origin=independent-derivation
\begin{solution}
\textbf{补写解答（重构新增）。}设 $S\subseteq X$，$X$ 完全正则。每个 $s\in S$ 的单点在 $X$ 中闭，故在 $S$ 中闭，$S$ 为 $T_1$。任给 $S$ 中闭集 $E$ 与点 $s\notin E$，由相对闭性存在 $X$ 中闭集 $F$ 使 $E=F\cap S$。因为 $s\in S$，必有 $s\notin F$。完全正则性给出 $f\in C(X,[0,1])$，满足 $f(s)=1$ 且 $f|_F=0$。连续限制 $f|_S$ 在 $s$ 取 $1$、在 $E$ 取 $0$，所以 $S$ 完全正则。$S=\varnothing$ 时所需点集分离条件真空成立，结论仍对。
\end{solution}
\begin{remark}
\textbf{条件说明（重构新增）。}相对拓扑使每个相对闭集能写成环境闭集与子空间的交，并使连续函数的限制仍连续。这两点是继承性的全部机制，不要求 $S$ 在 $X$ 中闭或开。
\end{remark}
% END SUPPLEMENT: id=v14-ch04-ex72-solution
% SOURCE p=161; printed=146; object=ch04-ex73-compactification-algebras; number=73; region=upper
\begin{exercise}\label{exercise:ch04-compactification-algebras}
设 $X$ 为完全正则空间。
\begin{definition*}[完全正则子代数]\phantomsection\label{definition:ch04-completely-regular-algebra}
若 $BC(X)$ 的子代数 $\mathcal A$ 满足：\textup{(i)} 它闭且包含常函数；\textup{(ii)} $\mathcal A\cap C(X,I)$ 分离点与闭集，则称 $\mathcal A$ \textbf{完全正则}。
\end{definition*}
\begin{enumerate}
\item\label{item:ch04-compactification-algebra-regular} 若 $(Y,e)$ 是 $X$ 的 Hausdorff 紧化，则 $\mathcal A_Y=\{f\circ e:f\in C(Y)\}$ 是 $BC(X)$ 的完全正则子代数。
\item\label{item:ch04-compactification-algebra-unique} 若 $(Y,e)$、$(Y',e')$ 是 $X$ 的 Hausdorff 紧化，且 $\mathcal A_Y=\mathcal A_{Y'}$，则存在同胚 $\phi:Y\to Y'$ 使 $\phi\circ e=e'$。（改用\cref{theorem:ch04-stone-cech-universal}的证明；该定理处理 $Y=\beta X$ 的情形。）
\item\label{item:ch04-compactification-generated-algebra} 若 $(Y,e)$ 是与 $\mathcal F\subseteq C(X,I)$ 关联的 $X$ 的紧化，则 $\mathcal A_Y$ 是 $BC(X)$ 中包含 $\mathcal F$ 的最小闭子代数。（使用\cref{exercise:ch04-cube-coordinate-algebra}。）
\item\label{item:ch04-compactification-algebra-correspondence} $X$ 的 Hausdorff 紧化与 $BC(X)$ 的完全正则子代数一一对应。
\end{enumerate}
\end{exercise}
% SUPPLEMENT: id=v14-ch04-ex73-solution; parent=ch04-ex73-compactification-algebras; kind=solution+source-correction+counterexample+motivation+hypotheses; origin=independent-derivation
\begin{remark}
\textbf{定义的动机与原题勘误（重构新增）。}一个紧化可以由“哪些有界连续函数能延拓到它”来编码。闭性反映一致极限仍能延拓，常函数总能延拓，分离点与闭集保证该代数足以恢复 $X$ 的拓扑。这里仍遵守本书 $BC(X)$ 为复函数代数、子代数不自动含单位的约定。

按此约定，原第三项遗漏了常函数 $1$，第四项还遗漏了对复共轭封闭的要求。两项原文已保留，下面分别给出反例，并证明正确版本；不把不成立的原陈述当作已证。第四项的紧化应按“存在与嵌入相容的同胚”取等价类，避免将同一紧化的不同命名计作不同对象。
\end{remark}
\begin{solution}
\textbf{前两项及修正版的完整解答（重构新增）。}
\begin{enumerate}
\item 限制映射 $R:C(Y)\to BC(X)$、$Rf=f\circ e$ 保持线性、乘法和常函数。因 $e(X)$ 稠密，$\|Rf\|_u=\|f\|_u$；其像 $\mathcal A_Y$ 由完备性为闭子代数，理由是像中任意收敛列的原像为 Cauchy 列，其极限仍在 $C(Y)$。

若 $E\subseteq X$ 闭且 $x\notin E$，则 $e(E)$ 在 $e(X)$ 中相对闭，可写 $e(E)=F\cap e(X)$，其中 $F$ 在 $Y$ 中闭。$e(x)\notin F$。紧 Hausdorff 空间 $Y$ 正规，故存在 $h\in C(Y,I)$ 使 $h(e(x))=1$、$h|_F=0$。于是 $h\circ e\in\mathcal A_Y\cap C(X,I)$，在 $x$ 为 $1$、在 $E$ 为 $0$，所需分离性质成立。因此 $\mathcal A_Y$ 完全正则。还可记录后面需要的事实：它自动对共轭封闭，因为 $\overline{h\circ e}=\overline h\circ e$。

\item 令 $\mathcal G=C(Y',I)$，$i':Y'\to I^{\mathcal G}$ 为关联嵌入，其像因 $Y'$ 紧而闭。对每个 $g\in\mathcal G$，$g\circ e'\in\mathcal A_{Y'}=\mathcal A_Y$，故存在 $h_g\in C(Y)$ 满足
\begin{align*}
h_g\circ e=g\circ e'.
\end{align*}
稠密性使 $h_g$ 唯一，并且使它在整个 $Y$ 上取值于 $I$：闭集 $h_g^{-1}(I)$ 包含 $e(X)$，所以是 $Y$。定义 $H:Y\to I^{\mathcal G}$ 为 $\pi_g\circ H=h_g$，则 $H$ 连续，且 $H\circ e=i'\circ e'$。对稠密集取闭包得
\begin{align*}
H(Y)\subseteq\overline{H(e(X))}
=\overline{i'(e'(X))}=i'(Y').
\end{align*}
因此 $\phi=(i')^{-1}\circ H:Y\to Y'$ 连续，且 $\phi\circ e=e'$。

对调 $Y,Y'$，同样构造连续 $\psi:Y'\to Y$，满足 $\psi\circ e'=e$。于是 $\psi\circ\phi$ 与 $Y$ 上恒等映射在稠密集 $e(X)$ 相等，因 $Y$ Hausdorff 而处处相等；同理 $\phi\circ\psi=\operatorname{id}_{Y'}$。故 $\phi$ 为同胚，且满足题目要求。若 $X$ 为空，两紧化也为空，唯一空映射直接给出结论。

\item \textbf{原陈述的反例。}取离散空间 $X=\mathbb N$，令 $\mathcal F$ 为所有有限支集的 $I$ 值函数。任意闭集 $E$ 与外部点 $n$ 可由 $\chi_{\{n\}}$ 分离，故这族函数确实给出关联紧化。由于离散 $\mathbb N$ 的紧子集恰为有限集，$\mathcal F=C_c(\mathbb N)\cap C(\mathbb N,I)$。前面的单点紧化验证表明，关联紧化为 $\mathbb N\cup\{\infty\}$，其中 $\infty$ 的邻域补集有限。

由 $\mathcal F$ 生成的不含单位要求的复代数恰为有限支集复序列：有限和与积仍有限支集，每个有限支集复序列又是有限个 $\chi_{\{n\}}$ 的复线性组合。其一致闭包为
\begin{align*}
c_0=\{(a_n):a_n\to0\}.
\end{align*}
确实，趋零序列由有限截断一致逼近；若一个序列是一致极限，则给定 $\varepsilon>0$ 先取一个误差小于 $\varepsilon$ 的有限支集近似，离开其有限支集后原序列模小于 $\varepsilon$，故趋零。

另一方面，函数延拓到 $\infty$ 连续，当且仅当序列收敛到其在 $\infty$ 的值，所以
\begin{align*}
\mathcal A_Y=c=\{(a_n):a_n\text{ 在 }\mathbb C\text{ 中收敛}\}.
\end{align*}
常序列 $1$ 属于 $c$ 却不属于 $c_0$，否定原等号。

\textbf{修正版：}$\mathcal A_Y$ 是包含 $\mathcal F$ \textbf{及常函数 $1$} 的最小闭子代数。为证明它，令 $\mathcal B$ 为 $C(Y)$ 中由 $1$ 及坐标限制 $\pi_f|_Y$（$f\in\mathcal F$）生成的复代数。它含 $1$；坐标均实值，故它对共轭封闭；$Y$ 为立方体子空间，不同点必有不同坐标，所以它分离点。复 Stone--Weierstrass 定理给 $\overline{\mathcal B}=C(Y)$。限制等距映射 $R$ 把生成元送到 $1$ 及 $\mathcal F$，故把 $\mathcal B$ 送到它们在 $BC(X)$ 中的生成代数。等距性与第一项中的闭像性质给
\begin{align*}
\mathcal A_Y=R(C(Y))=\overline{R(\mathcal B)},
\end{align*}
即所需最小闭子代数。这里用的是\cref{exercise:ch04-cube-coordinate-algebra}所用的同一逼近论证，直接在闭子空间 $Y$ 上实施，避免额外引入延拓步骤。

\item \textbf{原复代数陈述的反例。}令 $\mathbb T=\{z\in\mathbb C:|z|=1\}$，$\mathcal P$ 为复多项式在 $\mathbb T$ 上的一致闭包。它是含常数的闭复子代数，含坐标函数 $z$，却不含 $\overline z$；\cref{example:ch04-conjugate-circle-obstruction}已证明对任意多项式 $p$ 都有 $\|\overline z-p(z)\|_u\ge1$。

取 $X=\mathbb T\times\mathbb N$，其中 $\mathbb N$ 离散。各层 $\mathbb T\times\{n\}$ 都开闭，且为紧 Hausdorff 空间，因此 $X$ 为 LCH，特别完全正则。$BC(X)$ 的元素可写成序列 $(f_n)$，其中 $f_n\in C(\mathbb T)$ 且 $\sup_n\|f_n\|_u<\infty$；反向这样的序列在每个开层连续，从而在 $X$ 连续。定义
\begin{align*}
\mathcal A=\{(f_n)\in BC(X):
\text{存在 }h\in\mathcal P\text{，使 }\|f_n-h\|_u\to0\}.
\end{align*}
它含常函数，并且是复子代数：线性组合的尾部极限为对应线性组合；乘积的尾部极限为对应乘积，因为各收敛函数列一致有界，乘积差可用两个范数乘积估计，而 $\mathcal P$ 是代数。

验证闭性。若 $a^{(k)}=(f_n^{(k)})\in\mathcal A$ 在 $BC(X)$ 中趋于 $a=(f_n)$，记各尾部极限为 $h_k\in\mathcal P$。对每个 $k,l$，让 $n\to\infty$ 得
\begin{align*}
\|h_k-h_l\|_u\le\sup_n\|f_n^{(k)}-f_n^{(l)}\|_u
=\|a^{(k)}-a^{(l)}\|_u.
\end{align*}
所以 $(h_k)$ 是 Cauchy 列，趋于闭完备子空间 $\mathcal P$ 中某个 $h$。任给 $\varepsilon>0$，先选 $k$ 使 $\|a-a^{(k)}\|_u<\varepsilon/3$ 且 $\|h_k-h\|_u<\varepsilon/3$，再选 $N$ 使 $n\ge N$ 时 $\|f_n^{(k)}-h_k\|_u<\varepsilon/3$。三角不等式给 $\|f_n-h\|_u<\varepsilon$，故 $a\in\mathcal A$。因此 $\mathcal A$ 闭。

它还满足原定义的分离要求。设 $x=(z_0,n_0)\notin E$，$E$ 在 $X$ 中闭。截面 $E_{n_0}=\{z:(z,n_0)\in E\}$ 在 $\mathbb T$ 中闭，且不含 $z_0$。由 Urysohn 引理，选 $u\in C(\mathbb T,I)$，使 $u(z_0)=1$、$u|_{E_{n_0}}=0$。在第 $n_0$ 层取 $u$、其余层取零，得到 $f\in\mathcal A\cap C(X,I)$，因为它的尾部极限为 $0\in\mathcal P$；它在 $x$ 为 $1$，在 $E$ 为 $0$。

但是每一层都取 $z$ 的函数属于 $\mathcal A$，其共轭却不属于：共轭序列的唯一一致尾部极限为 $\overline z\notin\mathcal P$。因此 $\mathcal A$ 不对共轭封闭。第一项表明，每个真正紧化产生的 $\mathcal A_Y$ 都对共轭封闭，所以这个满足原定义的完全正则子代数不对应任何紧化，否定原第四项。

\textbf{修正版及证明。}$X$ 的 Hausdorff 紧化等价类，与 $BC(X)$ 中\textbf{对共轭封闭的完全正则子代数}一一对应。第一项将每个紧化送到此类代数，第二项说明代数相同则紧化等价。反过来，给定此类 $\mathcal A$，取 $\mathcal F=\mathcal A\cap C(X,I)$，形成关联紧化 $Y$。由修正后的第三项，$\mathcal A_Y$ 是 $\mathcal F$ 与 $1$ 生成的闭代数，故 $\mathcal A_Y\subseteq\mathcal A$。

任取 $f\in\mathcal A$。共轭封闭使
\begin{align*}
u=\frac{f+\overline f}{2},\qquad
v=\frac{f-\overline f}{2i}
\end{align*}
也属于 $\mathcal A$，它们为有界实函数。取 $M>\max(\|u\|_u,\|v\|_u)$，则 $(u+M)/(2M)$ 和 $(v+M)/(2M)$ 都属于 $\mathcal F$，因此 $u,v$ 及 $f=u+iv$ 属于由 $\mathcal F$ 与 $1$ 生成的代数。故 $\mathcal A\subseteq\mathcal A_Y$，两者相等。两个构造互逆到紧化等价关系，证明正确的一一对应。
\end{enumerate}
\end{solution}
\begin{remark}
\textbf{条件的作用（重构新增）。}第三项的常数解决“所有生成函数在新边界点同时为零”的障碍；第四项的共轭封闭使实部和虚部仍在代数内，才能由 $[0,1]$ 值函数恢复全部复函数。这是两个独立问题。对实函数代数，相应的实部、虚部障碍不存在；本题按书中的复函数约定处理，不能悄悄改变标量域来掩盖遗漏。
\end{remark}
% END SUPPLEMENT: id=v14-ch04-ex73-solution
% SOURCE p=161; printed=146; object=ch04-ex74-stone-cech-integers; number=74; region=middle
\begin{exercise}\label{exercise:ch04-stone-cech-integers}
把取离散拓扑的 $\mathbb N$ 视为其 Stone--Čech 紧化 $\beta\mathbb N$ 的子集。
\begin{enumerate}
\item\label{item:ch04-stone-cech-disjoint-closures} 若 $A,B$ 是 $\mathbb N$ 的不交子集，则它们在 $\beta\mathbb N$ 中的闭包也不交。（提示：$\chi_A\in C(\mathbb N,I)$。）
\item\label{item:ch04-stone-cech-no-nonconstant-sequence} $\mathbb N$ 中的序列，除非最终为常数，否则不会在 $\beta\mathbb N$ 中收敛。因此 $\beta\mathbb N$ 确实不是序列紧的。
\end{enumerate}
\end{exercise}
% SUPPLEMENT: id=v14-ch04-ex74-solution; parent=ch04-ex74-stone-cech-integers; kind=solution+motivation+hypotheses; origin=independent-derivation
\begin{solution}
\textbf{补写解答（重构新增）。}
\begin{enumerate}
\item 离散拓扑使 $\chi_A$ 连续，故它有连续延拓 $h:\beta\mathbb N\to I$。闭集 $h^{-1}(\{1\})$ 包含 $A$，所以也包含 $\overline A$；闭集 $h^{-1}(\{0\})$ 包含 $B$，所以也包含 $\overline B$。两个原像不交，得到 $\overline A\cap\overline B=\varnothing$。
\item 设 $(n_k)$ 在 $\beta\mathbb N$ 中收敛到 $p$。若它取无限多个不同值，可依次选取下标严格递增且值两两不同的子列 $(n_{k_j})$：每次已选有限个值，而后续不可能只取这些值，否则整列的值域有限。令
\begin{align*}
A=\{n_{k_{2j-1}}:j\ge1\},\qquad
B=\{n_{k_{2j}}:j\ge1\}.
\end{align*}
二者不交，奇、偶子列均收敛到 $p$，所以 $p\in\overline A\cap\overline B$，与第一项矛盾。

因此原序列只取有限多个值。若不最终为常数，则至少两个不同值各出现无限多次：若只有一个值出现无限次，其余有限多个值共只出现有限次，序列就会最终恒为该值。两个无限常值子列分别收敛到这两个不同点，又都必须收敛到 $p$，违背 Hausdorff 空间极限唯一性。故序列最终为常数。

最后，$1,2,3,\ldots$ 是 $\beta\mathbb N$ 中的一列，其任意子列仍取两两不同的自然数，不能最终为常数，因而没有收敛子列。所以紧空间 $\beta\mathbb N$ 不是序列紧空间。
\end{enumerate}
\end{solution}
\begin{remark}
\textbf{动机与条件说明（重构新增）。}离散性使任意集合的示性函数都连续，这比通常空间的连续函数分离能力强得多。闭包不交由这一函数一直保持的 $0/1$ 差别保证。第二项只讨论各项位于原 $\mathbb N$ 的序列；用这类序列已足以否定整个紧化的序列紧性，不能把结论误写成对 $\beta\mathbb N$ 内任意序列的未经证明的断言。
\end{remark}
% END SUPPLEMENT: id=v14-ch04-ex74-solution
% SOURCE p=161; printed=146; object=ch04-ex75-stone-cech-characters; number=75; region=middle-lower
\begin{exercise}\label{exercise:ch04-stone-cech-characters}
设 $X$ 为完全正则空间。从 $BC(X,\mathbb R)$ 到 $\mathbb R$ 的全部非零代数同态组成的集合 $M$，赋予逐点收敛拓扑后，同胚于 $\beta X$。（见\cref{exercise:ch04-recover-space-from-algebra}。从 Banach 代数理论的观点看，这是 $\beta X$ 的自然实现。）
\end{exercise}
% SUPPLEMENT: id=v14-ch04-ex75-solution; parent=ch04-ex75-stone-cech-characters; kind=solution+motivation+hypotheses; origin=independent-derivation
\begin{solution}
\textbf{补写解答（重构新增）。}连续延拓的存在、唯一与一致范数等距性表明，限制映射
\begin{align*}
R:C(\beta X,\mathbb R)\longrightarrow BC(X,\mathbb R),
\qquad Rg=g\circ e
\end{align*}
是实代数的等距同构。满射性由每个有界连续实函数均能延拓得到，实值性由稠密集上的实值与 $\mathbb R$ 在 $\mathbb C$ 中闭得到。

令 $M_{\beta X}$ 是 $C(\beta X,\mathbb R)$ 的非零实同态空间。映射
\begin{align*}
T:M\to M_{\beta X},\qquad T(\phi)=\phi\circ R
\end{align*}
为双射，逆映射为 $\psi\mapsto\psi\circ R^{-1}$；非零性由 $R$ 满射及单射保持。两个空间分别是 $\mathbb R^{BC(X,\mathbb R)}$ 与 $\mathbb R^{C(\beta X,\mathbb R)}$ 的子空间。对每个 $g\in C(\beta X,\mathbb R)$，$T(\phi)(g)=\phi(Rg)$ 是源空间的一个坐标，所以 $T$ 连续；对每个 $f\in BC(X,\mathbb R)$，$T^{-1}(\psi)(f)=\psi(R^{-1}f)$ 也为坐标，逆映射连续。因此 $T$ 为同胚。

\cref{exercise:ch04-recover-space-from-algebra}用于紧 Hausdorff 空间 $\beta X$，给 $\beta X\cong M_{\beta X}$；再与 $T^{-1}$ 复合即得所需同胚。具体地，$p\in\beta X$ 对应同态 $f\mapsto\widetilde f(p)$，其中 $\widetilde f$ 为唯一连续延拓。若 $X$ 为空，各空间为空，结论仍成立。
\end{solution}
\begin{remark}
\textbf{动机与条件说明（重构新增）。}原空间的点对应通常求值，紧化边界点对应“对全部有界连续函数一致兼容的求值规则”。同态不预设连续，紧空间的代数恢复定理已经证明它们自动为求值。同胚结论需要逐点收敛拓扑，不能只证明代数同态集合之间存在双射就结束。
\end{remark}
% END SUPPLEMENT: id=v14-ch04-ex75-solution
% SOURCE p=161; printed=146; object=ch04-ex76-countable-separating-family; number=76; region=lower
\begin{exercise}\label{exercise:ch04-countable-separating-family}
若 $X$ 正规且第二可数，则存在可数函数族 $\mathcal F\subseteq C(X,I)$ 分离点与闭集。（令 $\mathcal B$ 为拓扑的一个可数基。考虑满足 $\overline U\subseteq V$ 的点对 $(U,V)\in\mathcal B\times\mathcal B$，并使用 Urysohn 引理。）
\end{exercise}
% SUPPLEMENT: id=v14-ch04-ex76-solution; parent=ch04-ex76-countable-separating-family; kind=solution+motivation+hypotheses; origin=independent-derivation
\begin{solution}
\textbf{补写解答（重构新增）。}取可数基 $\mathcal B$。对每个满足 $\overline U\subseteq V$ 的 $(U,V)\in\mathcal B^2$，闭集 $\overline U$ 与 $X\setminus V$ 不交，故 Urysohn 引理给出 $f_{U,V}\in C(X,I)$，使它在 $\overline U$ 上为 $1$，在 $X\setminus V$ 上为 $0$。满足条件的点对至多可数，所以这些函数组成的族至多可数；加入常函数 $1$ 后仍可数且非空。

现设 $E$ 闭且 $x\notin E$。先取 $V\in\mathcal B$ 使 $x\in V\subseteq X\setminus E$。正规 $T_1$ 空间正则，所以存在开集 $W$，使 $x\in W\subseteq\overline W\subseteq V$。再取 $U\in\mathcal B$ 使 $x\in U\subseteq W$。于是
\begin{align*}
\overline U\subseteq\overline W\subseteq V,
\end{align*}
该点对已列入构造。相应 $f_{U,V}(x)=1$，且 $E\subseteq X\setminus V$，故 $f_{U,V}|_E=0$。因此所选可数族分离点与闭集。
\end{solution}
\begin{remark}
\textbf{动机与逐条件说明（重构新增）。}不能为每个点和闭集任取一个函数再宣称所得族可数，因为这些对象可能不可数。这里先用可数基的有序对作索引，再证明每次分离任务都由某个索引覆盖。闭包包含 $\overline U\subseteq V$ 留下缓冲区域，保证两侧闭集不交；正规性制造函数，正则性保证足够多这样的基对，而本书的 $T_1$ 约定保证正规蕴含正则。
\end{remark}
% END SUPPLEMENT: id=v14-ch04-ex76-solution
% SOURCE p=161; printed=146; object=ch04-ex77-countable-product-metric; number=77; region=lower
\begin{exercise}\label{exercise:ch04-countable-product-metric}
设 $\{(X_n,\rho_n)\}_{n=1}^{\infty}$ 为可数族度量空间，且各度量取值于 $[0,1]$。后一限制总可满足，见\subcref{exercise:ch04-bounded-function-metrics}{item:ch04-bounded-metric}。令 $X=\prod_{n=1}^{\infty}X_n$。对 $x,y\in X$，写 $x=(x_1,x_2,\ldots)$、$y=(y_1,y_2,\ldots)$，定义
\begin{align*}
\rho(x,y)=\sum_{n=1}^{\infty}2^{-n}\rho_n(x_n,y_n).
\end{align*}
则 $\rho$ 是 $X$ 上给出乘积拓扑的度量。
\end{exercise}
% SUPPLEMENT: id=v14-ch04-ex77-solution; parent=ch04-ex77-countable-product-metric; kind=solution+motivation+hypotheses; origin=independent-derivation
\begin{solution}
\textbf{补写解答（重构新增）。}若某个因子为空，$X$ 为空，唯一空度量诱导唯一拓扑，结论直接成立。以下设 $X$ 非空。每项非负且至多为 $2^{-n}$，所以级数收敛到 $[0,1]$ 内的有限数。对称性由各项得到；若 $\rho(x,y)=0$，所有非负项都为零，所以每个 $x_n=y_n$，即 $x=y$；反向显然。对每个 $n$ 用三角不等式，先对前 $N$ 项求和再令 $N\to\infty$，得到 $\rho(x,z)\le\rho(x,y)+\rho(y,z)$，故 $\rho$ 为度量。

证明两个拓扑相同。先取 $x$ 的一个乘积基本邻域 $U$，它只限制有限坐标集 $F$。对每个 $n\in F$，取 $\varepsilon_n>0$，使 $\rho_n$-球 $B_{\rho_n}(x_n,\varepsilon_n)$ 包含于该坐标的指定开邻域。若 $F$ 非空，令 $\delta=\min_{n\in F}2^{-n}\varepsilon_n>0$。当 $\rho(x,y)<\delta$ 时，对 $n\in F$ 有
\begin{align*}
2^{-n}\rho_n(x_n,y_n)\le\rho(x,y)<\delta
\le2^{-n}\varepsilon_n,
\end{align*}
所以 $y\in U$。故 $B_\rho(x,\delta)\subseteq U$。若 $F$ 为空，则 $U=X$，任取正 $\delta$ 即可。这证明每个乘积开集为 $\rho$-开集。

反向任给 $\varepsilon>0$，取 $N$ 使 $2^{-N}<\varepsilon/2$，并考虑乘积基本邻域
\begin{align*}
V=\{y\in X:\rho_n(x_n,y_n)<\varepsilon/2
\text{ 对 }1\le n\le N\}.
\end{align*}
对 $y\in V$，有
\begin{align*}
\rho(x,y)
&=\sum_{n=1}^N2^{-n}\rho_n(x_n,y_n)
+\sum_{n>N}2^{-n}\rho_n(x_n,y_n)\\
&<\frac\varepsilon2\sum_{n=1}^N2^{-n}+2^{-N}<\varepsilon.
\end{align*}
故 $V\subseteq B_\rho(x,\varepsilon)$。任意 $\rho$-开集的每个点都有包含于该开集的乘积基本邻域，因此它也是乘积开集，两拓扑一致。
\end{solution}
\begin{remark}
\textbf{动机与逐条件说明（重构新增）。}可数性让所有坐标能排成序列，并分为有限前段与统一很小的尾段。上界 $1$ 保证尾和对任意点一致受控；权重 $2^{-n}$ 可换成任意严格正、可求和的权重，正性保留每个坐标的辨别能力。各因子的度量可先用前述有界化方法替换为等价有界度量，因此上界不限制它们的拓扑。有限个因子的情形可在后面补单点因子；这也覆盖度量化定理中分离族有限的情况。
\end{remark}
% END SUPPLEMENT: id=v14-ch04-ex77-solution
\end{document}
